\documentclass[11pt,a4paper]{article}
\usepackage[utf8]{inputenc}
\usepackage[T1]{fontenc}
\usepackage{geometry}
\usepackage{hyperref}
\usepackage{setspace}

\geometry{margin=1in}
\onehalfspacing

\begin{document}

\begin{flushright}
    Krzysztof Żuchowski \\
    Independent Researcher \\
    Fractal Information Theory Project \\
    Email: \href{mailto:krzysztof.zuch@gmail.com}{krzysztof.zuch@gmail.com} \\
    Date: \today
\end{flushright}

\vspace{1cm}

\begin{flushleft}
    To: \\
    The Editorial Office \\
    \textit{Foundations of Physics}
\end{flushleft}

\vspace{0.5cm}

\noindent Dear Editor,

\bigskip

\noindent I hereby submit the manuscript entitled

\medskip

\noindent \textbf{``Global Fractal Temporal Correlations in Gravitational-Wave Detector Noise: Cross-Detector, Cross-Epoch, and Null-Model Validation (QW–1660 Phase 16)''}

\medskip

\noindent for consideration for publication in \textit{Foundations of Physics}.

\bigskip

\noindent This work presents a comprehensive, multi-stage statistical investigation of long-range temporal correlations in publicly available gravitational-wave interferometer strain data from the LIGO Hanford (H1), LIGO Livingston (L1), and Virgo (V1) detectors. The analysis encompasses Phase 16 of the QW–1660 research program (versions v5–v24) and is based exclusively on open data and fully reproducible methods.

\bigskip

\noindent The central objective of the study is to test, under increasingly stringent controls, whether the standard assumption of stationary, Gaussian, and locally correlated detector noise fully accounts for the observed statistical structure of long-duration strain data. Using a consistent multi-scale rescaled-range (R/S) estimator, we identify stable anti-persistent fractal scaling in individual detectors across multiple observational epochs (O3a, O3b, and early O4). More importantly, we observe non-zero cross-Hurst exponents between geographically separated detectors.

\bigskip

\noindent A key component of the work is rigorous null-model validation. We demonstrate that the observed cross-detector correlations are strongly suppressed by temporal shuffling, which destroys long-range memory, while remaining partially preserved under phase-randomized surrogates that conserve the power spectral density. This combination of results excludes explanations based solely on shared spectral features, coincident transient contamination, or simple stationary noise models.

\bigskip

\noindent The manuscript does not propose modifications to General Relativity and does not rely on speculative assumptions. All findings are formulated strictly at the level of statistical structure and are fully compatible with standard local spacetime dynamics. The potential implications are discussed conservatively, in terms of emergent or global statistical organization of spacetime fluctuations, a topic that lies squarely within the conceptual scope of \textit{Foundations of Physics}.

\bigskip

\noindent The primary contribution of this work is both methodological and empirical. It provides a transparent, reproducible framework for probing long-range temporal correlations in gravitational-wave data that complements existing spectral and transient-based analyses, and it establishes a set of falsifiable statistical observations that warrant further investigation.

\bigskip

\noindent The manuscript has not been published previously and is not under consideration for publication elsewhere. All data and code required to reproduce the results are publicly available and explicitly referenced in the manuscript.

\bigskip

\noindent Thank you for your consideration. I would be grateful for the opportunity to have this work reviewed for publication in \textit{Foundations of Physics}.

\bigskip

\noindent Sincerely,

\vspace{1cm}

\noindent Krzysztof Żuchowski \\
Independent Researcher \\
Fractal Information Theory Project

\end{document}
